%% SCRAP: architecture/03-architecture/adaptive-systems/window-inference-redesign
%% SOURCE: docs/working/architecture/03-architecture/adaptive-systems/window-inference-redesign.md
%% STATUS: WORKING
%% FITS: dev-guide/ch-physics
%% EDITORIAL: lifted — prose rewritten to press voice

\section{Window-Width Inference: A Statistically Valid Redesign}

This section sets out a redesign, dated 2025-11-19 and approved for
implementation at the time of writing, of the window-width inference algorithm
used by the adaptive runtime. The original \texttt{find\_variance\_inflection()}
measured the wrong quantity --- prefix variance --- and the redesign replaces it
with a hypothesis test for variance stability across disjoint windows.

\subsection{Why the Original Algorithm Was Invalid}

The original algorithm scanned candidate window sizes and, for each, computed the
variance of the trajectory prefix from the start through that size, stopping when
the change fell below one percent of the full variance. Four statistical defects
follow. The samples are not independent: a larger prefix wholly contains every
smaller one, so the variance estimates are correlated and biased. The data are
not stationary: execution heat decays roughly as $\textit{heat}(t) = h_0\,
e^{-s t}$, so early elements are hot and late elements cold, and both mean and
variance drift with position. The decision is confounded: a flattening variance
cannot be distinguished between the window being wide enough for the workload (the
goal) and the prefix simply accumulating low-heat samples that dilute variance (an
artifact). And the threshold is arbitrary: one percent is a magic number with no
statistical justification and no accompanying confidence measure.

The failure is concrete. On a trajectory that decays through hot, warm, and cold
phases, the algorithm halts not where the workload's true pattern width lies but
where it runs out of hot data and crosses into the low-variance cold region --- a
stopping point with nothing to do with the correct window width.

\subsection{The Redesign: Variance-Stability Testing}

The redesign replaces prefix variance with stability testing over disjoint
windows. For each candidate size $N$, the trajectory is divided into
non-overlapping chunks of size $N$, the variance of each chunk is computed
independently, and a statistical test asks whether those chunk variances are
equal. The smallest $N$ at which the variances are statistically similar is the
sufficient window width.

\subsection{Levene's Test}

The chosen instrument is Levene's test for equality of variance, evaluated in
\Qtype{} fixed point. The null hypothesis is that all chunk variances are equal.
The test statistic is

\begin{equation}
W = \frac{(K-1)\sum_i n_i (z_i - \bar{z})^2}
         {\sum_i \sum_j (z_{ij} - z_i)^2},
\end{equation}

where $K$ is the number of chunks, $n_i$ the size of chunk $i$,
$z_{ij} = |x_{ij} - \mathrm{median}_i|$ the absolute deviation from the chunk
median, $z_i$ the chunk mean of those deviations, and $\bar{z}$ their overall
mean. When $W$ exceeds the critical value at $\alpha = 0.05$ the variances differ
significantly and a larger window is needed; when $W$ falls at or below it the
variances are similar enough and the window is sufficient. The redesign uses a
conservative critical value near 6.5 to accommodate integer arithmetic.

Three edge cases are handled explicitly: a trajectory shorter than the candidate
window falls back to the minimum window; reaching the maximum window without
passing the test returns the maximum; and a candidate that yields fewer than
three chunks is skipped, since reliable testing needs at least three.

\subsection{Implementation Outline}

The Levene statistic is added as a helper over an array of per-chunk variances,
supported by median and mean routines in fixed point. The inference routine then
scans candidate sizes in steps of 64 from the minimum to the trajectory-bounded
maximum, computes per-chunk variances, applies the test, and returns the first
size that passes. The inference output structure is extended with diagnostics ---
the Levene statistic, the number of chunks tested, a pass flag, and the minimum
sufficient window --- so the result is inspectable from external analysis tooling.

\subsection{Validation}

Unit tests anchor the test against known inputs: four identically-variance chunks
must pass with a statistic far below the critical value, four monotonically
increasing variances must fail with a statistic well above it, and a synthetic
trajectory with a known pattern width must return a window inside the valid
range. Regression testing confirms the existing suite still passes, and the old
and new algorithms are run side by side on the same trajectories for comparison.

\subsection{Expected Outcome}

The redesign trades an ad-hoc heuristic for a grounded hypothesis test. Where the
old algorithm violated independence, confounded decay with sufficiency, leaned on
a magic threshold, and could vary run to run, the new one uses disjoint windows,
tests within-chunk stability, reports a defensible confidence measure, and is
deterministic. The expectation is that window-width estimates stabilize across
runs and that the resulting optimization becomes reproducible and defensible. The
assessed risk is medium complexity with no breaking interface change and roughly
ten to twenty percent more computation, still under one percent of total VM
overhead.

\subsection{References}

\begin{itemize}
  \item Levene, H. (1960). ``Robust tests for equality of variances.'' In
        \textit{Contributions to Probability and Statistics}, ed. I. Olkin et
        al. Stanford University Press.
  \item Brown, M. B., \& Forsythe, A. B. (1974). ``Robust tests for the equality
        of variances.'' \textit{Journal of the American Statistical
        Association}, 69(346), 364--367.
  \item NIST/SEMATECH e-Handbook of Statistical Methods, section on Levene's
        test.
\end{itemize}
